\paragraph{Quantitative normalized-penalty identity.}
The constant-penalty criterion admits a simpler
quantitative form that requires no assumption on
the sign of $U_n$. We record it separately because
it gives a sharper criterion and a finite-penalty
error estimate independent of RH.

\begin{theorem}[Unconditional quantitative penalty limit]
\label{thm:quantpenalty}
For fixed $n\ge1$ and $0<T\le\tfrac12$, let
$U_n(s,x)=L_{n,\,1/2-s}(x)$ and
$k_n=(n+1)(2n+1)$. Define
\[
M_{n,T}=k_n\sup_{\substack{0\le s\le T\\x\in\mathbb R}}
 |U_{n+1}(s,x)|<\infty,
\]
and the nonnegative heat-smoothed negative part
\[
\mathcal N_n(T,x)=\int_0^T
 \bigl(P_{T-s}(-U_n(s,\cdot))_+\bigr)(x)\,ds.
\]
For every $C>0$ and every real $x$,
\begin{equation}\label{eq:quantpenalty}
\boxed{
\left|\frac{\mathcal E_{n,C}(T,x)}{C}
      -\mathcal N_n(T,x)\right|
 \le\frac{T M_{n,T}}{C}.}
\end{equation}
The convergence as $C\to\infty$ is uniform
in $x$ for every fixed $n,T$.
No Laguerre positivity hypothesis is used.
\end{theorem}
\begin{proof}
Write $v=k_nU_{n+1}(s,y)$ and $u=U_n(s,y)$.
The positive-part map $a\mapsto a_+$
is $1$-Lipschitz, so
\[
\left|[v/C-u]_+-[-u]_+\right|
\le |v|/C\le M_{n,T}/C.
\]
Since $C>0$,
$[v/C-u]_+=[v-Cu]_+/C$.
The Gaussian semigroup $P_h$ preserves
positivity, constants, and the supremum
norm. Apply it to this pointwise difference
and integrate over $0\le s\le T$,
which gives \eqref{eq:quantpenalty}.
Finally, the Fourier integral for $H_t$
and its derivatives has a uniform
majorant on the compact heat-time slab:
for each integer $j\ge0$,
\[
\sup_{\substack{0\le t\le1/2\\x\in\mathbb R}}
 |H_t^{(j)}(x)|
\le \int_0^\infty
u^j e^{u^2/2}\Phi(u)\,du<\infty.
\]
The finite derivative-polynomial definition
of $L_{n+1,t}$ therefore implies
$M_{n,T}<\infty$.
\end{proof}

\begin{corollary}[A bounded-penalty criterion for RH]
\label{cor:boundedpenaltyRH}
For the genuine theta heat family, RH
is equivalent to each of the following
statements:
\begin{enumerate}
\item For every $n\ge1$ and real $x$,
$\sup_{C\ge1}\mathcal E_{n,C}(\tfrac12,x)<\infty$.
\item For every $n\ge1$ and real $x$,
$\mathcal E_{n,C}(\tfrac12,x)=o(C)$
as $C\to\infty$.
\end{enumerate}
Both conditions are statements about
\emph{all} Laguerre orders and all real
Fourier arguments. Neither is established
by the tests in this paper.
\end{corollary}
\begin{proof}
The normalized limit
\eqref{eq:quantpenalty} identifies (2)
with $\mathcal N_n(\tfrac12,x)=0$.
Because the Gaussian heat kernel is
strictly positive for positive heat time,
and $U_n$ is continuous, the latter
is equivalent to $U_n(s,y)\ge0$
on the entire slab. Indeed, any strictly
negative space--time value extends to
a rectangle of positive measure and
contributes strictly positively to
$\mathcal N_n(\tfrac12,x)$.
The already proved all-order
Laguerre--P\'olya characterization
then gives RH.
Condition (1) implies (2).
Conversely, RH makes all $U_n$
nonnegative on the heat slab by the
de Bruijn--Newman theorem, whence
$E_{n,C}\le[k_nU_{n+1}]_+\le M_{n,T}$
for all $C$ and
$\mathcal E_{n,C}(T,x)\le TM_{n,T}$,
establishing (1).
\end{proof}

\begin{corollary}[A finite-$C$ negative-mass bound]
\label{cor:negative-mass}
For every $C>0$,
\[
0\le\mathcal N_n(T,x)
\le\frac{\mathcal E_{n,C}(T,x)
 +T M_{n,T}}{C}.
\]
Conversely, if
$\mathcal N_n(T,x)\ge\delta>0$, then
\[
\mathcal E_{n,C}(T,x)
\ge C\delta-TM_{n,T},
\]
giving a quantitative linear-growth
certificate for an off-axis obstruction.
\end{corollary}

\paragraph{What is and is not gained.}
The earlier unscaled condition
$\mathcal E_{n,C}\to0$ has been
replaced, for the purpose of a
sufficient RH criterion, by the
strictly weaker demand of
\emph{boundedness in $C$}.
This is a genuine simplification of
the analytic target, but the bound
has not been proved for the theta
correlation measure. It remains an
RH-equivalent condition, not an
independent proof of RH.
One cannot infer it from finite
Laguerre orders, finite Fourier
heights or from the central-frequency
positivity results already obtained.

For a sharp elementary illustration
take $U(s,x)=1-2s$, $v=2$,
and $T=1$. If $C>2$, direct
integration gives
\[
\mathcal E_C(1,x)=C/4+1+1/C,
\qquad
\mathcal N(1,x)=1/4.
\]
Thus
$\mathcal E_C/C-\mathcal N=1/C+1/C^2$:
the error is $O(C^{-1})$, while the
unscaled penalty grows linearly.
The exact finite comparison and
rational model checks are recorded in
\path{XiKernel/PenaltyQuantitative.lean}
and
\path{certificates/xi_penalty_quantitative_exact.py}.
The Lean statements formalize the
finite algebra, not the heat-semigroup
integrals or the RH equivalence.
